\documentclass[conference]{IEEEtran}

\usepackage{amsmath}
\usepackage{graphicx}
\usepackage{booktabs}
\usepackage{hyperref}

\title{Tracklet Repair and Conservative Merging for More Stable Single-Camera Tracking}

\author{
    \IEEEauthorblockN{Author Name}
    \IEEEauthorblockA{
        Deep Learning Lab\\
        University Name\\
        Email Address
    }
}

\begin{document}

\maketitle

\begin{abstract}
This paper studies a post-processing approach for improving single-camera
tracking consistency inside a multi-camera multi-target tracking pipeline. The
method focuses on tracklet analysis, short-gap repair, and conservative merging
of fragmented tracklets. This is a placeholder abstract and should be updated
after experiments are completed.
\end{abstract}

\section{Introduction}

Single-camera tracking errors can propagate into multi-camera association. This
project investigates whether simple and conservative post-processing can improve
tracklet stability before cross-camera matching.

\section{Related Work}

TODO: Summarize relevant work on multi-object tracking, tracklet association,
and multi-camera tracking.

\section{Method}

TODO: Describe the baseline tracker output, tracklet analysis, gap repair rules,
and conservative merge criteria.

\section{Experiments}

TODO: Describe the dataset, metrics, implementation details, and compared
settings.

\section{Results}

TODO: Add quantitative tables and qualitative examples.

\section{Discussion}

TODO: Discuss limitations, failure cases, and trade-offs between stability and
identity correctness.

\section{Conclusion}

TODO: Summarize the main findings after experiments are complete.

\bibliographystyle{IEEEtran}
\bibliography{references}

\end{document}
